\subsection{Extension of derivations: the case of rings}

Let K be a commutative ring, A a commutative K-algebra and \(\mathbf{x} = (x_i)_{i \in I}\) a family of elements of A. Further, let \(\Delta\) be a derivation of K into an A-module M, in other words (III, p.~553) a Z-linear mapping of K into M satisfying the relation \(\Delta(cc') = c \cdot \Delta(c') + c' \cdot \Delta(c)\) for c, \(c'\) in K. For each \(i \in I\) let \(D_i\) be the partial derivation with respect to \(X_i\) in the polynomial ring \(K[X_i]_{i \in I}\) ; this is the unique derivation of that ring into itself which is zero on K and on \(X_j\) for \(j \in I - \{i\}\) , and takes the value 1 on \(X_i\) (IV, p.~6). For every polynomial \(f = \sum_{\alpha \in N^{(1)}} c_\alpha \cdot X^\alpha\) in \(K[X_i]_{i \in I}\) we denote by \(f^\Delta(\mathbf{x})\) the element \(\sum_{\alpha \in N^{(1)}} x^\alpha \cdot \Delta(c_\alpha)\) of

\begin{enumerate}
\def\labelenumi{\Alph{enumi}.}
\setcounter{enumi}{12}
\tightlist
\item
\end{enumerate}

PROPOSITION 1. --- Suppose that the family \(\mathbf{x} = (x_i)_{i \in I}\) generates the K-algebra A. Let \((m_i)_{i \in I}\) be a family of elements of M and \((f_\lambda)_{\lambda \in \Lambda}\) a family of polynomials generating the ideal \(\mathfrak{a}\) of all polynomials \(f \in \mathbf{K}[\mathbf{X}_i]_{i \in I}\) such that \(f(\mathbf{x}) = 0\) . For a derivation D of A into M to exist such that D(c.1) = \(\Delta(c)\) for all c \(\in\) K and D(x\_i) = m\_i for all i \(\in\) I it is necessary and sufficient that

\[
f _ {\lambda} ^ {\Delta} (\mathbf {x}) + \sum_ {i \in I} D _ {i} f _ {\lambda} (\mathbf {x}). m _ {i} = 0 \quad \text { for   all } \quad \lambda \in \Lambda .\tag{1}
\]

If this is so, the derivation D is unique and satisfies

\[
\mathrm{D} (f (\mathbf {x})) = f ^ {\Delta} (\mathbf {x}) + \sum_ {i \in I} \mathrm{D} _ {i} f (\mathbf {x}). m _ {i} \quad \text { for   all } \quad f \in \mathrm{K} [ \mathrm{X} _ {i} ] _ {i \in I}.\tag{2}
\]

Put \(\mathrm{E} = \mathrm{K}[\mathrm{X}_i]_{i \in I}\) and write \(\varphi\) for the K-homomorphism of E onto A which maps \(\mathbf{X}_i\) to \(x_i\) for all \(i \in I\) ; we thus have \(\varphi(f) = f(\mathbf{x})\) for all \(f \in E\) . We consider M as an E-module by means of the homomorphism \(\varphi: \mathrm{E} \to \mathrm{A}\) and define a mapping \(D'\) of E into M by \(D'(f) = f^{\Delta}(\mathbf{x}) + \sum_{i \in I} D_i f(\mathbf{x}) \cdot m_i\) (note that the family \((D_i f)_{i \in I}\) has finite support for each \(f \in E\) ). It is clear that \(D'\) is the unique derivation of E into M extending \(\Delta\) and mapping \(X_i\) to \(m_i\) for each \(i \in I\) .

Let D be a derivation of A into M such that \(\mathrm{D}(c,1)=\Delta(c)\) for all \(c\in K\) and \(\mathrm{D}(x_{i})=m_{i}\) for all \(i\in I\) . Then \(D\circ\varphi\) is a derivation of E into M extending \(\Delta\) and mapping \(X_{i}\) to \(m_{i}\) for all \(i\in I\) . We thus have \(D\circ\varphi=D'\) , that is, the relation (2) holds. This proves the uniqueness of D; moreover (1) is a consequence of \(f_{\lambda}(\mathbf{x})=0\) and (2).

Conversely assume that (1) holds; in other words, we have \(\mathbf{D}'(f_{\lambda}) = 0\) for all \(\lambda \in \Lambda\) . Let \(f \in \mathfrak{a}\) ; there exists a family with finite support \((q_{\lambda})_{\lambda \in \Lambda}\) in E such that \(f = \sum_{\lambda \in \Lambda} q_{\lambda} \cdot f_{\lambda}\) . We have

\[
\mathrm{D} ^ {\prime} (f) = \sum_ {\lambda \in \Lambda} \left[ f _ {\lambda} (\mathbf {x}) \cdot \mathrm{D} ^ {\prime} (q _ {\lambda}) + q _ {\lambda} (\mathbf {x}) \cdot \mathrm{D} ^ {\prime} (f _ {\lambda}) \right]
\]

whence \(D'(f) = 0\) because \(f_{\lambda}(\mathbf{x})\) and \(D'(f_{\lambda})\) are zero for all \(\lambda \in \Lambda\) . Since \(D'\) vanishes on \(\mathfrak{a}\), there exists a Z-linear mapping D of A into M such that \(D' = D \circ \varphi\) ; clearly D is the required derivation of A into M.

